% Generado por scripts/hoja_de_curacion.py. NO EDITAR A MANO.
\documentclass[10pt,a4paper]{article}
\usepackage[utf8]{inputenc}
\usepackage[T1]{fontenc}
\usepackage[spanish,es-nodecimaldot,es-noquoting]{babel}
\usepackage[a4paper,margin=2.1cm]{geometry}
\usepackage{lmodern}
\usepackage{booktabs}
\usepackage{array}
\usepackage{amssymb}
\usepackage{xcolor}
\usepackage{mdframed}
\usepackage{microtype}
\usepackage{titlesec}
\usepackage{fancyhdr}
\usepackage[hidelinks]{hyperref}

% LOS IDENTIFICADORES DE MATHLIB TRAEN UNICODE, y pdfLaTeX no lo sabe leer.
% El primer intento murio con «Unicode character U+2080 not set up» por un
% `min_bi` con subindice. Se declara el RANGO, no el caracter que fallo hoy:
% esta hoja se regenera, y manana los modulos pendientes seran otros.
\DeclareUnicodeCharacter{2080}{\ensuremath{_0}}
\DeclareUnicodeCharacter{2081}{\ensuremath{_1}}
\DeclareUnicodeCharacter{2082}{\ensuremath{_2}}
\DeclareUnicodeCharacter{2083}{\ensuremath{_3}}
\DeclareUnicodeCharacter{2084}{\ensuremath{_4}}
\DeclareUnicodeCharacter{2085}{\ensuremath{_5}}
\DeclareUnicodeCharacter{2086}{\ensuremath{_6}}
\DeclareUnicodeCharacter{2087}{\ensuremath{_7}}
\DeclareUnicodeCharacter{2088}{\ensuremath{_8}}
\DeclareUnicodeCharacter{2089}{\ensuremath{_9}}
\DeclareUnicodeCharacter{1D62}{\ensuremath{_i}}
\DeclareUnicodeCharacter{2C7C}{\ensuremath{_j}}
\DeclareUnicodeCharacter{207F}{\ensuremath{^n}}
\DeclareUnicodeCharacter{03B1}{\ensuremath{\alpha}}
\DeclareUnicodeCharacter{03B2}{\ensuremath{\beta}}
\DeclareUnicodeCharacter{03B3}{\ensuremath{\gamma}}
\DeclareUnicodeCharacter{03B4}{\ensuremath{\delta}}
\DeclareUnicodeCharacter{03BC}{\ensuremath{\mu}}
\DeclareUnicodeCharacter{03C3}{\ensuremath{\sigma}}
\DeclareUnicodeCharacter{03C6}{\ensuremath{\varphi}}
\DeclareUnicodeCharacter{2115}{\ensuremath{\mathbb{N}}}
\DeclareUnicodeCharacter{2124}{\ensuremath{\mathbb{Z}}}
\DeclareUnicodeCharacter{211A}{\ensuremath{\mathbb{Q}}}
\DeclareUnicodeCharacter{211D}{\ensuremath{\mathbb{R}}}
\DeclareUnicodeCharacter{2102}{\ensuremath{\mathbb{C}}}
\DeclareUnicodeCharacter{1D55C}{\ensuremath{\Bbbk}}
\DeclareUnicodeCharacter{2192}{\ensuremath{\rightarrow}}
\DeclareUnicodeCharacter{2200}{\ensuremath{\forall}}
\DeclareUnicodeCharacter{2203}{\ensuremath{\exists}}
% Los simbolos de CATEGORIAS, que salen en cuanto la hoja habla de alias:
% `SimplexCategoryGenRel ≌ SimplexCategory` tumbo la compilacion. Se declara
% la familia entera y no el que fallo hoy, que es la leccion de U+2080.
\DeclareUnicodeCharacter{224C}{\ensuremath{\backsimeq}}
\DeclareUnicodeCharacter{2243}{\ensuremath{\simeq}}
\DeclareUnicodeCharacter{2245}{\ensuremath{\cong}}
\DeclareUnicodeCharacter{27F9}{\ensuremath{\Longrightarrow}}
\DeclareUnicodeCharacter{27F6}{\ensuremath{\longrightarrow}}
\DeclareUnicodeCharacter{2964}{\ensuremath{\rightarrowtail}}
\DeclareUnicodeCharacter{2045}{\ensuremath{[\![}}
\DeclareUnicodeCharacter{2046}{\ensuremath{]\!]}}
\DeclareUnicodeCharacter{1D7ED}{\ensuremath{\mathbf{1}}}

\definecolor{acento}{RGB}{70,60,140}
\definecolor{suave}{RGB}{120,120,130}
\definecolor{fondo}{RGB}{244,243,248}
\definecolor{aviso}{RGB}{150,90,20}

\newcommand{\Mathlib}{\textsc{Mathlib}}

\titleformat{\section}{\large\bfseries\color{acento}}{}{0pt}{}
\titlespacing{\section}{0pt}{16pt}{6pt}
\titleformat{\subsection}{\normalsize\bfseries}{}{0pt}{}
\titlespacing{\subsection}{0pt}{10pt}{4pt}

\newmdenv[backgroundcolor=fondo,linewidth=0pt,skipabove=8pt,skipbelow=8pt,
          innerleftmargin=10pt,innerrightmargin=10pt,
          innertopmargin=8pt,innerbottommargin=8pt]{marca}
\newmdenv[linecolor=aviso,linewidth=1.2pt,topline=false,bottomline=false,
          rightline=false,skipabove=6pt,skipbelow=6pt,
          innerleftmargin=8pt,innertopmargin=4pt,
          innerbottommargin=4pt]{aviso}

\pagestyle{fancy}\fancyhf{}
\renewcommand{\headrulewidth}{0.3pt}
\fancyhead[L]{\small\color{suave}Metamatemático · curación pendiente}
\fancyhead[R]{\small\color{suave}\thepage}

\begin{document}
\begin{center}
{\LARGE\bfseries\color{acento} Curación pendiente}\par\smallskip
{\color{suave}\small Los 41 módulos que el grafo no cubre, con todo lo
verificable ya resuelto.}\par
{\color{suave}\footnotesize Generado por \texttt{scripts/hoja\_de\_curacion.py}
· se regenera, no se edita a mano}
\end{center}
\vspace{4pt}

\section*{Dos grupos, dos apuestas}

El banco de enunciados que fallan está \textbf{agotado}: los 41 módulos que
destapó se curaron y los 10 que quedaban están decididos. Esta hoja sale de
mirar directamente qué partes de \Mathlib{} no cubre el grafo, y trae
\textbf{24} módulos en dos grupos que conviene no
mezclar.\par\medskip
\textbf{Grupo 1 · seguir la tanda} —12 módulos. La
curación anterior metió el primer nodo de varias ramas y se paró ahí.
\texttt{Algebra.Lie} tiene \textbf{1\,228 teoremas en 50 ficheros y el grafo
toca uno}: un álgebra de Lie sin subálgebras, sin ideales y sin pesos. Son
rincones —poco volumen— pero el barrio se mide limpio.\par\medskip
\textbf{Grupo 2 · el núcleo} —12 módulos.
El mismo criterio sobre \Mathlib{} entero da \textbf{161 ramas abiertas y
finas}, y las mayores no son las de la tanda: \texttt{Algebra.Order} con
5\,140 teoremas y \textbf{1 módulo cubierto de 208}, \texttt{Algebra.Group}
con 4\,102 y 4 de 148, \texttt{Analysis.SpecialFunctions} con 3\,906 y 1 de
98. Ahí caen las consultas de un alumno, y explica el \textbf{5,0\,\%} de
precisión sobre \Mathlib{} entero.\par\medskip
\textbf{Una pregunta extra para el grupo 2.} Ahí el grafo YA tiene nodo de
cabecera —\texttt{group-theory}, \texttt{real-analysis}— y lo que falta son
los módulos finos de debajo. Puede que la respuesta correcta no sea un nodo
nuevo sino \emph{más nombres en el que ya hay}, que es más barato. Si es el
caso, escríbelo en vez de la marca.

\begin{marca}
\textbf{1 · La marca.} Es lo que sostiene el grafo, y no la decide nada automático:\par\smallskip
\begin{tabular}{@{}llll@{}}\toprule
\texttt{C} & una categoría & $\to$ & \textbf{vértice}\\
\texttt{S} & una subcategoría plena & $\to$ & \textbf{vértice}, y la inclusión es arista\\
\texttt{F} & un funtor o clase de flechas & $\to$ & \textbf{arista} \emph{en este ambiente}\\
\texttt{O} & un objeto individual & $\to$ & vértice degenerado\\
\texttt{T} & ni objetos ni flechas & $\to$ & \textbf{fuera}\\
\bottomrule\end{tabular}\par\smallskip
\begin{sloppypar}
\texttt{F} \textbf{no} dice «esto no puede ser un vértice nunca». Eso es falso, y está demostrado que lo es en \texttt{FlechasComoObjetos.lean}: las flechas de $C$ son exactamente los objetos de la categoría de flechas $\mathrm{Arrow}\,C$, y los funtores son los objetos de la categoría de funtores $[C,D]$. Lo que dice es que \textbf{en este grafo} —cuyos objetos son conceptos y cuyas flechas son dependencias— la etiqueta nombra una flecha.\end{sloppypar}\smallskip
\texttt{homology} es \texttt{F}: aquí no es una colección que se pueda colimitar, es el funtor \emph{a lo largo del cual} se colimita. \texttt{prime-factorization} es \texttt{T}: es un teorema, no un objeto.\par\medskip
\textbf{2 · El padre.} De qué concepto es especialización. La flecha va del general al específico.\par\medskip
\textbf{3 · Cuáles nombres se quedan.} Que exista no lo hace correcto: \texttt{Nat.Prime} es la noción de primo, \texttt{Nat.minFac} no lo es.
\end{marca}

\noindent Lo que \textbf{no} hay que decidir, porque ya está verificado: si el
nombre existe, en qué módulo vive, cuál es el canónico —las citas— y el DAG de
\texttt{import}s. Los identificadores de abajo están \textbf{leídos de su
declaración}, no deducidos de la ruta: de 447 deducidos así, 95 no existían.
\par\medskip

\begin{center}\small\begin{tabular}{@{}lr@{\quad}lr@{\quad}lr@{}}\toprule
Algebra & 7 & Computability & 4 & Data & 4\\
Analysis & 3 & Dynamics & 2 & SetTheory & 2\\
NumberTheory & 1 & Order & 1\\
\bottomrule\end{tabular}\end{center}
\section{Algebra\hfill{\normalsize\normalfont área \texttt{algebra}}}
\begin{sloppypar}\noindent\small\textit{Padres candidatos ya en el grafo:} \texttt{abelian-groups}, \texttt{bases-and-dimension}, \texttt{bilinear-forms}, \texttt{canonical-forms}, \texttt{character-modules}, \texttt{character-theory}, \texttt{commutative-algebra}, \texttt{derivations}, \texttt{derived-category}, \texttt{different-ideal}\end{sloppypar}
\subsection*{\texttt{Algebra.Group.Basic}}
\noindent\small Sin sustantivos propios: sólo aporta teoremas. El grafo aporta \emph{sustantivos} —de sus 176 identificadores ninguno es un teorema—, así que probablemente no le toca.\par\medskip
\subsection*{\texttt{Algebra.Group.Pointwise.Finset.Basic}}
\noindent\begin{minipage}{0.56\linewidth}\small\begin{tabular}{@{}>{\raggedright\arraybackslash}p{0.58\linewidth}ll@{}}\toprule
identificador & tipo & citas\\\midrule
\texttt{Finset.\allowbreak{}coeMonoidHom} & def & 5\\
\texttt{Finset.\allowbreak{}singletonMulHom} & def & 4\\
\texttt{Finset.\allowbreak{}zpow} & def & 4\\
\texttt{Finset.\allowbreak{}singleton\allowbreak{}Monoid\allowbreak{}Hom} & def & 4\\
\bottomrule\end{tabular}\end{minipage}\hfill\begin{minipage}{0.40\linewidth}\small$\square$~marca \texttt{C S F O T}\par\smallskip$\square$~padre: \dotfill\par\smallskip$\square$~se quedan: \dotfill\par\end{minipage}\par\medskip
\subsection*{\texttt{Algebra.Lie.Nilpotent}}
\noindent\begin{minipage}{0.56\linewidth}\small\begin{tabular}{@{}>{\raggedright\arraybackslash}p{0.58\linewidth}ll@{}}\toprule
identificador & tipo & citas\\\midrule
\texttt{LieModule.\allowbreak{}IsNilpotent} & class & 251\\
\texttt{LieModule.\allowbreak{}lower\allowbreak{}Central\allowbreak{}Series} & def & 57\\
\texttt{LieRing.\allowbreak{}IsNilpotent} & abbrev & 10\\
\texttt{LieAlgebra.\allowbreak{}max\allowbreak{}Nilpotent\allowbreak{}Ideal} & def & 5\\
\bottomrule\end{tabular}\end{minipage}\hfill\begin{minipage}{0.40\linewidth}\small$\square$~marca \texttt{C S F O T}\par\smallskip$\square$~padre: \dotfill\par\smallskip$\square$~se quedan: \dotfill\par\end{minipage}\par\medskip
\subsection*{\texttt{Algebra.Lie.Subalgebra}}
\noindent\begin{minipage}{0.56\linewidth}\small\begin{tabular}{@{}>{\raggedright\arraybackslash}p{0.58\linewidth}ll@{}}\toprule
identificador & tipo & citas\\\midrule
\texttt{LieHom.\allowbreak{}range} & def & 2184\\
\texttt{LieSubalgebra.\allowbreak{}comap} & def & 1132\\
\texttt{LieSubalgebra.\allowbreak{}inclusion} & def & 169\\
\texttt{LieSubalgebra} & structure & 77\\
\bottomrule\end{tabular}\end{minipage}\hfill\begin{minipage}{0.40\linewidth}\small$\square$~marca \texttt{C S F O T}\par\smallskip$\square$~padre: \dotfill\par\smallskip$\square$~se quedan: \dotfill\par\end{minipage}\par\medskip
\subsection*{\texttt{Algebra.Lie.Submodule}}
\noindent\begin{minipage}{0.56\linewidth}\small\begin{tabular}{@{}>{\raggedright\arraybackslash}p{0.58\linewidth}ll@{}}\toprule
identificador & tipo & citas\\\midrule
\texttt{LieModuleHom.\allowbreak{}range} & def & 2184\\
\texttt{LieSubmodule.\allowbreak{}comap} & def & 1132\\
\texttt{LieSubmodule.\allowbreak{}inclusion} & def & 169\\
\texttt{LieSubmodule} & structure & 107\\
\bottomrule\end{tabular}\end{minipage}\hfill\begin{minipage}{0.40\linewidth}\small$\square$~marca \texttt{C S F O T}\par\smallskip$\square$~padre: \dotfill\par\smallskip$\square$~se quedan: \dotfill\par\end{minipage}\par\medskip
\subsection*{\texttt{Algebra.Order.Monoid.Unbundled.Basic}}
\noindent\begin{minipage}{0.56\linewidth}\small\begin{tabular}{@{}>{\raggedright\arraybackslash}p{0.58\linewidth}ll@{}}\toprule
identificador & tipo & citas\\\midrule
\texttt{MulLECancellable} & def & 20\\
\texttt{Contravariant.\allowbreak{}to\allowbreak{}Left\allowbreak{}Cancel\allowbreak{}Semigroup} & def & 0\\
\texttt{Contravariant.\allowbreak{}to\allowbreak{}Right\allowbreak{}Cancel\allowbreak{}Semigroup} & def & 0\\
\bottomrule\end{tabular}\end{minipage}\hfill\begin{minipage}{0.40\linewidth}\small$\square$~marca \texttt{C S F O T}\par\smallskip$\square$~padre: \dotfill\par\smallskip$\square$~se quedan: \dotfill\par\end{minipage}\par\medskip
\subsection*{\texttt{Algebra.Order.ToIntervalMod}}
\noindent\begin{minipage}{0.56\linewidth}\small\begin{tabular}{@{}>{\raggedright\arraybackslash}p{0.58\linewidth}ll@{}}\toprule
identificador & tipo & citas\\\midrule
\texttt{toIocMod} & def & 106\\
\texttt{toIcoMod} & def & 99\\
\texttt{toIcoDiv} & def & 66\\
\texttt{toIocDiv} & def & 66\\
\bottomrule\end{tabular}\end{minipage}\hfill\begin{minipage}{0.40\linewidth}\small$\square$~marca \texttt{C S F O T}\par\smallskip$\square$~padre: \dotfill\par\smallskip$\square$~se quedan: \dotfill\par\end{minipage}\par\medskip
\section{Computability\hfill{\normalsize\normalfont área \texttt{computation}}}
\begin{sloppypar}\noindent\small\textit{Padres candidatos ya en el grafo:} \texttt{algorithm-analysis}, \texttt{computability-theory}, \texttt{computational-complexity}, \texttt{finite-automata}, \texttt{formal-languages}, \texttt{formal-verification}, \texttt{lambda-calculus}, \texttt{np-completeness}, \texttt{partial-functions}, \texttt{partial-recursive-functions}\end{sloppypar}
\subsection*{\texttt{Computability.Primrec.List}}
\noindent\begin{minipage}{0.56\linewidth}\small\begin{tabular}{@{}>{\raggedright\arraybackslash}p{0.58\linewidth}ll@{}}\toprule
identificador & tipo & citas\\\midrule
\texttt{Nat.\allowbreak{}Primrec'} & inductive & 23\\
\texttt{Nat.\allowbreak{}Primrec'.\allowbreak{}Vec} & def & 0\\
\bottomrule\end{tabular}\end{minipage}\hfill\begin{minipage}{0.40\linewidth}\small$\square$~marca \texttt{C S F O T}\par\smallskip$\square$~padre: \dotfill\par\smallskip$\square$~se quedan: \dotfill\par\end{minipage}\par\medskip
\subsection*{\texttt{Computability.Reduce}}
\noindent\begin{minipage}{0.56\linewidth}\small\begin{tabular}{@{}>{\raggedright\arraybackslash}p{0.58\linewidth}ll@{}}\toprule
identificador & tipo & citas\\\midrule
\texttt{toNat} & def & 93\\
\texttt{ManyOneDegree.\allowbreak{}liftOn} & abbrev & 17\\
\texttt{ManyOneEquiv} & def & 16\\
\texttt{ManyOneDegree} & def & 12\\
\bottomrule\end{tabular}\end{minipage}\hfill\begin{minipage}{0.40\linewidth}\small$\square$~marca \texttt{C S F O T}\par\smallskip$\square$~padre: \dotfill\par\smallskip$\square$~se quedan: \dotfill\par\end{minipage}\par\medskip
\subsection*{\texttt{Computability.TMToPartrec}}
\noindent\begin{minipage}{0.56\linewidth}\small\begin{tabular}{@{}>{\raggedright\arraybackslash}p{0.58\linewidth}ll@{}}\toprule
identificador & tipo & citas\\\midrule
\texttt{Turing.\allowbreak{}PartrecToTM2.\allowbreak{}head} & def & 153\\
\texttt{Turing.\allowbreak{}PartrecToTM2.\allowbreak{}init} & def & 49\\
\texttt{Turing.\allowbreak{}PartrecToTM2.\allowbreak{}Supports} & def & 22\\
\texttt{Turing.\allowbreak{}PartrecToTM2.\allowbreak{}trNormal} & def & 17\\
\bottomrule\end{tabular}\end{minipage}\hfill\begin{minipage}{0.40\linewidth}\small$\square$~marca \texttt{C S F O T}\par\smallskip$\square$~padre: \dotfill\par\smallskip$\square$~se quedan: \dotfill\par\end{minipage}\par\medskip
\subsection*{\texttt{Computability.Tape}}
\noindent\begin{minipage}{0.56\linewidth}\small\begin{tabular}{@{}>{\raggedright\arraybackslash}p{0.58\linewidth}ll@{}}\toprule
identificador & tipo & citas\\\midrule
\texttt{Turing.\allowbreak{}proj} & def & 183\\
\texttt{Turing.\allowbreak{}ListBlank} & def & 45\\
\texttt{Turing.\allowbreak{}Tape} & structure & 16\\
\texttt{Turing.\allowbreak{}Tape.\allowbreak{}mk'} & def & 14\\
\bottomrule\end{tabular}\end{minipage}\hfill\begin{minipage}{0.40\linewidth}\small$\square$~marca \texttt{C S F O T}\par\smallskip$\square$~padre: \dotfill\par\smallskip$\square$~se quedan: \dotfill\par\end{minipage}\par\medskip
\section{Data\hfill{\normalsize\normalfont \textbf{sin área en el grafo}}}
\begin{aviso}\Mathlib{} no organiza esta rama como un área del grafo. Puede que estos conceptos sean \texttt{T} y no toquen.\end{aviso}
\subsection*{\texttt{Data.Finset.Lattice.Fold}}
\noindent\begin{minipage}{0.56\linewidth}\small\begin{tabular}{@{}>{\raggedright\arraybackslash}p{0.58\linewidth}ll@{}}\toprule
identificador & tipo & citas\\\midrule
\texttt{Finset.\allowbreak{}sup'} & def & 31\\
\texttt{Finset.\allowbreak{}inf'} & def & 25\\
\texttt{Finset.\allowbreak{}sup} & def & 7\\
\texttt{Finset.\allowbreak{}inf} & def & 0\\
\bottomrule\end{tabular}\end{minipage}\hfill\begin{minipage}{0.40\linewidth}\small$\square$~marca \texttt{C S F O T}\par\smallskip$\square$~padre: \dotfill\par\smallskip$\square$~se quedan: \dotfill\par\end{minipage}\par\medskip
\subsection*{\texttt{Data.Set.Function}}
\noindent\begin{minipage}{0.56\linewidth}\small\begin{tabular}{@{}>{\raggedright\arraybackslash}p{0.58\linewidth}ll@{}}\toprule
identificador & tipo & citas\\\midrule
\texttt{Function.\allowbreak{}invFunOn} & def & 17\\
\bottomrule\end{tabular}\end{minipage}\hfill\begin{minipage}{0.40\linewidth}\small$\square$~marca \texttt{C S F O T}\par\smallskip$\square$~padre: \dotfill\par\smallskip$\square$~se quedan: \dotfill\par\end{minipage}\par\medskip
\subsection*{\texttt{Data.Set.Image}}
\noindent\small Sin sustantivos propios: sólo aporta teoremas. El grafo aporta \emph{sustantivos} —de sus 176 identificadores ninguno es un teorema—, así que probablemente no le toca.\par\medskip
\subsection*{\texttt{Data.Set.Lattice}}
\noindent\begin{minipage}{0.56\linewidth}\small\begin{tabular}{@{}>{\raggedright\arraybackslash}p{0.58\linewidth}ll@{}}\toprule
identificador & tipo & citas\\\midrule
\texttt{Set.\allowbreak{}sigmaToiUnion} & def & 3\\
\texttt{Set.\allowbreak{}union\allowbreak{}Eq\allowbreak{}Sigma\allowbreak{}Of\allowbreak{}Disjoint} & def & 2\\
\texttt{Set.\allowbreak{}sigmaEquiv} & def & 1\\
\texttt{Set.\allowbreak{}sUnionPowersetGI} & def & 0\\
\bottomrule\end{tabular}\end{minipage}\hfill\begin{minipage}{0.40\linewidth}\small$\square$~marca \texttt{C S F O T}\par\smallskip$\square$~padre: \dotfill\par\smallskip$\square$~se quedan: \dotfill\par\end{minipage}\par\medskip
\section{Analysis\hfill{\normalsize\normalfont área \texttt{analysis}}}
\begin{sloppypar}\noindent\small\textit{Padres candidatos ya en el grafo:} \texttt{banach-spaces}, \texttt{complex-analysis}, \texttt{conformal-maps}, \texttt{contour-integration}, \texttt{differentiation}, \texttt{functional-analysis}, \texttt{harmonic-analysis}, \texttt{hilbert-spaces}, \texttt{holder-continuity}, \texttt{holomorphic-functions}\end{sloppypar}
\subsection*{\texttt{Analysis.Normed.Group.Basic}}
\noindent\begin{minipage}{0.56\linewidth}\small\begin{tabular}{@{}>{\raggedright\arraybackslash}p{0.58\linewidth}ll@{}}\toprule
identificador & tipo & citas\\\midrule
\texttt{norm\allowbreak{}Group\allowbreak{}Seminorm} & def & 1\\
\texttt{normGroupNorm} & def & 1\\
\texttt{SeminormedGroup.\allowbreak{}induced} & abbrev & 0\\
\texttt{Seminormed\allowbreak{}Comm\allowbreak{}Group.\allowbreak{}induced} & abbrev & 0\\
\bottomrule\end{tabular}\end{minipage}\hfill\begin{minipage}{0.40\linewidth}\small$\square$~marca \texttt{C S F O T}\par\smallskip$\square$~padre: \dotfill\par\smallskip$\square$~se quedan: \dotfill\par\end{minipage}\par\medskip
\subsection*{\texttt{Analysis.SpecialFunctions.Pow.NNReal}}
\noindent\begin{minipage}{0.56\linewidth}\small\begin{tabular}{@{}>{\raggedright\arraybackslash}p{0.58\linewidth}ll@{}}\toprule
identificador & tipo & citas\\\midrule
\texttt{NNReal.\allowbreak{}rpow} & def & 5\\
\texttt{ENNReal.\allowbreak{}rpow} & def & 5\\
\texttt{NNReal.\allowbreak{}orderIsoRpow} & def & 2\\
\texttt{ENNReal.\allowbreak{}orderIsoRpow} & def & 2\\
\bottomrule\end{tabular}\end{minipage}\hfill\begin{minipage}{0.40\linewidth}\small$\square$~marca \texttt{C S F O T}\par\smallskip$\square$~padre: \dotfill\par\smallskip$\square$~se quedan: \dotfill\par\end{minipage}\par\medskip
\subsection*{\texttt{Analysis.SpecialFunctions.Pow.Real}}
\noindent\begin{minipage}{0.56\linewidth}\small\begin{tabular}{@{}>{\raggedright\arraybackslash}p{0.58\linewidth}ll@{}}\toprule
identificador & tipo & citas\\\midrule
\texttt{Real.\allowbreak{}rpow} & def & 5\\
\bottomrule\end{tabular}\end{minipage}\hfill\begin{minipage}{0.40\linewidth}\small$\square$~marca \texttt{C S F O T}\par\smallskip$\square$~padre: \dotfill\par\smallskip$\square$~se quedan: \dotfill\par\end{minipage}\par\medskip
\section{Dynamics\hfill{\normalsize\normalfont área \texttt{probability}}}
\begin{sloppypar}\noindent\small\textit{Padres candidatos ya en el grafo:} \texttt{brownian-motion}, \texttt{conditional-expectation}, \texttt{ergodic-theory}, \texttt{flows}, \texttt{limit-theorems}, \texttt{markov-chains}, \texttt{martingale-theory}, \texttt{probability-theory}, \texttt{random-variables}, \texttt{rotation-number}\end{sloppypar}
\subsection*{\texttt{Dynamics.PeriodicPts.Defs}}
\noindent\begin{minipage}{0.56\linewidth}\small\begin{tabular}{@{}>{\raggedright\arraybackslash}p{0.58\linewidth}ll@{}}\toprule
identificador & tipo & citas\\\midrule
\texttt{Function.\allowbreak{}minimalPeriod} & def & 57\\
\texttt{Function.\allowbreak{}IsPeriodicPt} & def & 49\\
\texttt{Function.\allowbreak{}periodicPts} & def & 29\\
\texttt{MulAction.\allowbreak{}period} & def & 23\\
\bottomrule\end{tabular}\end{minipage}\hfill\begin{minipage}{0.40\linewidth}\small$\square$~marca \texttt{C S F O T}\par\smallskip$\square$~padre: \dotfill\par\smallskip$\square$~se quedan: \dotfill\par\end{minipage}\par\medskip
\subsection*{\texttt{Dynamics.TopologicalEntropy.CoverEntropy}}
\noindent\begin{minipage}{0.56\linewidth}\small\begin{tabular}{@{}>{\raggedright\arraybackslash}p{0.58\linewidth}ll@{}}\toprule
identificador & tipo & citas\\\midrule
\texttt{Dynamics.\allowbreak{}coverMincard} & def & 22\\
\texttt{Dynamics.\allowbreak{}coverEntropy} & def & 22\\
\texttt{Dynamics.\allowbreak{}IsDynCoverOf} & def & 20\\
\texttt{Dynamics.\allowbreak{}cover\allowbreak{}Entropy\allowbreak{}Entourage} & def & 18\\
\bottomrule\end{tabular}\end{minipage}\hfill\begin{minipage}{0.40\linewidth}\small$\square$~marca \texttt{C S F O T}\par\smallskip$\square$~padre: \dotfill\par\smallskip$\square$~se quedan: \dotfill\par\end{minipage}\par\medskip
\section{SetTheory\hfill{\normalsize\normalfont área \texttt{set-theory}}}
\begin{sloppypar}\noindent\small\textit{Padres candidatos ya en el grafo:} \texttt{cardinal-arithmetic}, \texttt{descriptive-set-theory}, \texttt{forcing}, \texttt{hereditarily-finite-sets}, \texttt{large-cardinals}, \texttt{zfc-classes}, \texttt{zfc-sets}\end{sloppypar}
\subsection*{\texttt{SetTheory.ZFC.Ordinal}}
\noindent\begin{minipage}{0.56\linewidth}\small\begin{tabular}{@{}>{\raggedright\arraybackslash}p{0.58\linewidth}ll@{}}\toprule
identificador & tipo & citas\\\midrule
\texttt{ZFSet.\allowbreak{}IsOrdinal} & structure & 9\\
\texttt{ZFSet.\allowbreak{}IsTransitive} & def & 8\\
\texttt{Ordinal.\allowbreak{}toZFSet} & def & 8\\
\texttt{Ordinal.\allowbreak{}toPSet} & def & 0\\
\bottomrule\end{tabular}\end{minipage}\hfill\begin{minipage}{0.40\linewidth}\small$\square$~marca \texttt{C S F O T}\par\smallskip$\square$~padre: \dotfill\par\smallskip$\square$~se quedan: \dotfill\par\end{minipage}\par\medskip
\subsection*{\texttt{SetTheory.ZFC.PSet}}
\noindent\begin{minipage}{0.56\linewidth}\small\begin{tabular}{@{}>{\raggedright\arraybackslash}p{0.58\linewidth}ll@{}}\toprule
identificador & tipo & citas\\\midrule
\texttt{PSet.\allowbreak{}Nonempty} & def & 2513\\
\texttt{PSet.\allowbreak{}insert} & def & 922\\
\texttt{PSet.\allowbreak{}Equiv} & def & 518\\
\texttt{PSet.\allowbreak{}image} & def & 415\\
\bottomrule\end{tabular}\end{minipage}\hfill\begin{minipage}{0.40\linewidth}\small$\square$~marca \texttt{C S F O T}\par\smallskip$\square$~padre: \dotfill\par\smallskip$\square$~se quedan: \dotfill\par\end{minipage}\par\medskip
\section{NumberTheory\hfill{\normalsize\normalfont área \texttt{number-theory}}}
\begin{sloppypar}\noindent\small\textit{Padres candidatos ya en el grafo:} \texttt{algebraic-number-theory}, \texttt{analytic-number-theory}, \texttt{arithmetic-geometry}, \texttt{diophantine-equations}, \texttt{divisibility-gcd}, \texttt{elementary-number-theory}, \texttt{ideal-class-group}, \texttt{infinite-places}, \texttt{modular-arithmetic}, \texttt{modular-forms}\end{sloppypar}
\subsection*{\texttt{NumberTheory.ModularForms.QExpansion}}
\noindent\begin{minipage}{0.56\linewidth}\small\begin{tabular}{@{}>{\raggedright\arraybackslash}p{0.58\linewidth}ll@{}}\toprule
identificador & tipo & citas\\\midrule
\texttt{Slash\allowbreak{}Invariant\allowbreak{}Form\allowbreak{}Class.\allowbreak{}cuspFunction} & def & 39\\
\texttt{ModularFormClass.\allowbreak{}qExpansion} & def & 24\\
\texttt{UpperHalfPlane.\allowbreak{}valueAtInfty} & def & 5\\
\texttt{ModularFormClass.\allowbreak{}q\allowbreak{}Expansion\allowbreak{}Formal\allowbreak{}Multilinear\allowbreak{}Series} & def & 4\\
\bottomrule\end{tabular}\end{minipage}\hfill\begin{minipage}{0.40\linewidth}\small$\square$~marca \texttt{C S F O T}\par\smallskip$\square$~padre: \dotfill\par\smallskip$\square$~se quedan: \dotfill\par\end{minipage}\par\medskip
\section{Order\hfill{\normalsize\normalfont \textbf{sin área en el grafo}}}
\begin{aviso}\Mathlib{} no organiza esta rama como un área del grafo. Puede que estos conceptos sean \texttt{T} y no toquen.\end{aviso}
\subsection*{\texttt{Order.Filter.Map}}
\noindent\begin{minipage}{0.56\linewidth}\small\begin{tabular}{@{}>{\raggedright\arraybackslash}p{0.58\linewidth}ll@{}}\toprule
identificador & tipo & citas\\\midrule
\texttt{Filter.\allowbreak{}kernMap} & def & 6\\
\texttt{Filter.\allowbreak{}monad} & def & 0\\
\bottomrule\end{tabular}\end{minipage}\hfill\begin{minipage}{0.40\linewidth}\small$\square$~marca \texttt{C S F O T}\par\smallskip$\square$~padre: \dotfill\par\smallskip$\square$~se quedan: \dotfill\par\end{minipage}\par\medskip
\section*{Antes de dar por buena una tanda}

\begin{marca}
\noindent\texttt{python -m scripts.recuperacion\_contra\_proofnet}\par\smallskip
Precisión y cobertura contra 371 formalizaciones de oro, con su modelo nulo y
sin gastar API.\par\medskip
\textbf{Baseline hoy: 22,8\,\% / 18,0\,\% contra 1,45\,\% / 3,3\,\% —
15,7$\times$.} Si la precisión baja, esa tanda \textbf{no entra}.\par\medskip
Ya pasó: ofrecer los sustantivos de los nodos generados bajaba de 14,0\,\% a
11,5\,\%. Añadir vocabulario tiene coste.\par\medskip
\textbf{Y la regla tiene letra pequeña.} La primera tanda de 31 nodos bajaba
la precisión a 19,9\,\% sin mover la cobertura, y la causa no era ningún
veredicto equivocado —los 47 nombres los acepta \texttt{\#check}— sino que el
emparejador tokeniza el \textbf{id y el nombre} de cada nodo:
\texttt{different-\allowbreak{}ideal} aportaba el token \texttt{different},
que sale en media biblioteca, y con eso gastaba una de las dos plazas del
prompt. La puerta que lo arregla —\emph{sólo se ocupa plaza con una keyword
declarada}— subió la línea base de 21,3\,\% a 22,8\,\%. Si una tanda baja la
precisión, mira primero si sus nodos entran en plaza por su nombre.
\end{marca}

\end{document}